\section{Alternating Values: a Sum Bounded Through Three Lemmas}\label{sec:brs2}

\subsection{The program}
The test \texttt{sum/brs2}: the array alternates between $2$ and $0$, the sum is kept in the cell $0$ of a
one-element array, and the assertion is an upper bound.
\begin{lstlisting}[language=C]
for (i = 0; i < N; i++) { if (i % 2 == 0) a[i] = 2; else a[i] = 0; }
for (i = 0; i < N; i++) { if (i == 0) sum[0] = 0; else sum[0] = sum[0] + a[i]; }
__VERIFIER_assert(sum[0] <= 2 * N);
\end{lstlisting}

\subsection{The trace equation system}
Two loops in sequence; the second one's trace of \texttt{sum} has the final array $a'$ of the first in its
update. With $N > 0$ from the harness,
\begin{align}
  K &= \cnt(0, 1, n) \qquad a' = \cur(A) \qquad i' = \cur(I) \label{eq:brs0}\\
  A &= \tr{a_0}{\bl{K}, \bl{I}}{a, k, i \mapsto \upd{i \bmod 2 = 0}{a[i := 2]}\ \ \upd{i \bmod 2 \ne 0}{a[i := 0]}} \label{eq:brs1}\\
  I &= \tr{0}{\bl{K}}{i, k \mapsto \upd{i \bmod 2 = 0}{i + 1}\ \ \upd{i \bmod 2 \ne 0}{i + 1}} \label{eq:brs2}\\
  &\al{\bl{K}, \bl{A}, \bl{I}}{k, a, i \mapsto i < N} \qquad \neg(i' < N) \label{eq:brs3}\\
  K' &= \cnt(0, 1, n') \qquad s' = \cur(S) \qquad i'' = \cur(I') \nonumber\\
  S &= \tr{s_0}{\bl{K'}, \bl{I'}}{s, k, i \mapsto \upd{i = 0}{s[0 := 0]}\ \ \upd{i \ne 0}{s[0 := s[0] + a'[i]]}} \label{eq:brs4}\\
  I' &= \tr{0}{\bl{K'}}{i, k \mapsto i + 1} \qquad \al{\bl{K'}, \bl{S}, \bl{I'}}{k, s, i \mapsto i < N} \qquad \neg(i'' < N) \nonumber\\
  &\Longrightarrow\ s'[0] \le 2 N \nonumber
\end{align}
(the \texttt{rem} of C is the remainder of the prelude, which the simplifier turns into $\bmod$ under
$0 \le i$). The accumulator is not a scalar but a cell of an array, so \eqref{eq:brs4} is a trace of stores
at the constant key $0$, and the array it reads is the current value of another trace of stores.

\subsection{The lemmas}
Three lemmas of the database cooperate, in the order in which the goal needs them.
\begin{align*}
  &\lemma{stores\_index\_cases\_def}:\\
  &\quad \tr{a_0}{\bl{\cnt(0, 1, n)}, \bl{\cnt(0, 1, n)}}{a, k, i \mapsto \upd{P(k)}{a[k := c]}\ \ \upd{\neg P(k)}{a[k := d]}}\\
  &\qquad = \mathsf{stores}(a_0, \bl{\cnt(0, 1, n)}, \mapc{\bl{\cnt(0, 1, n)}}{k \mapsto \ite{P(k)}{c}{d}})\\[4pt]
  &\lemma{select\_current\_trace\_reset\_add}:\\
  &\quad \cur(\tr{a_0}{\bl{\cnt(0, 1, n)}, \bl{\cnt(0, 1, n)}}{a, k, i \mapsto \upd{k = 0}{a[j := c]}\ \ \upd{0 < k}{a[j := a[j] + G(k)]}})[j]\\
  &\qquad = \ite{n = 0}{a_0[j]}{c + \mathsf{sum}(\mapc{\bl{\cnt(0, 1, n)}}{k \mapsto \ite{0 < k}{G(k)}{0}})}\\[4pt]
  &\lemma{sum\_cases\_below}:\quad 0 \le e\ \wedge\ \intf(n) \cdot c \le e\ \wedge\ \intf(n) \cdot d \le e\\
  &\quad\Longrightarrow\ \mathsf{sum}(\mapc{\bl{\cnt(0, 1, n)}}{k \mapsto \ite{0 < k}{(\ite{P(k)}{c}{d})}{0}}) \le e\ \Longleftrightarrow\ \mathsf{true}
\end{align*}
The first turns a loop that stores at each position a constant chosen by cases on the \emph{index} into
stores of a map; the second reads a sum kept in one cell $j$ of an array as the scalar sum, with the
right-hand side of \lemma{current\_trace\_reset\_add} for a scalar accumulator, so that the database continues
from there; the third bounds a sum of $n$ terms that are each $0$, $c$ or $d$, with its left-hand side a
comparison whose constants bind the bound $e$ (the general \lemma{sum\_below\_prefix} of the database does not
apply, since its bound occurs only in a side condition, which would have to be a fact).

\subsection{The proof}
Of the 62 rule applications, those that carry the proof:
\begin{enumerate}[nosep]
\item \textbf{The counters.} \lemma{trace\_merge\_cases\_ne} joins the two cases of \eqref{eq:brs2}, both
  $i + 1$, and \lemma{counter\_def} makes $I = \cnt(0, 1, n)$, likewise $I' = \cnt(0, 1, n')$. In the
  context of the iteration $k$ the simplifier reads $i$ as $k$ (\lemma{earlier\_counter}), so the constraint
  \eqref{eq:brs3} is $\al{\bl{\cnt(0, 1, n)}}{k \mapsto k < N}$, which \lemma{invs\_counter\_upper\_nonnegative}
  resolves to $\intf(n) \le N$; with the exit $N \le i' = \intf(n)$ (\lemma{current\_counter}) this gives
  $\intf(n) = N$, and for the second loop $\intf(n') = N$, hence $n' = n$ (\lemma{integer\_injective}).
\item \textbf{The first loop is stores.} In the update of \eqref{eq:brs1}, $i \bmod 2 = 0$ is $k \bmod 2 = 0$
  and the guards are $k \bmod 2 = 0$ and $0 < k \bmod 2$; \lemma{stores\_index\_cases\_def} matches with
  $P := \lambda k.\ k \bmod 2 = 0$, $c := 2$, $d := 0$ (the higher-order $P$ by abstraction of the guard), so
  $a' = \cur(\mathsf{stores}(a_0, \bl{\cnt(0, 1, n)}, \mapc{\bl{\cnt(0, 1, n)}}{k \mapsto \ite{k \bmod 2 = 0}{2}{0}}))$.
\item \textbf{The read in the second loop.} The update of \eqref{eq:brs4} reads $a'[i]$, i.e.\ $a'[k]$;
  \lemma{current\_stores} makes it $\ite{\mathsf{occurs}(k, \bl{\cnt(0, 1, n)})}{\mathsf{earlier}(\mathsf{indexof}(k, \bl{\cnt(0, 1, n)}), \mapc{\ldots}{\ldots})}{a_0[k]}$,
  \lemma{occurs\_counter\_up} the test $0 \le k < \intf(n)$, which the context of the iteration decides (it
  has $0 \le k < n' = n$), \lemma{indexof\_counter\_up} the index $\natf(k)$, and the read of the map at $k$ is
  $\ite{k \bmod 2 = 0}{2}{0}$. The second loop is now
  $\tr{s_0}{\ldots}{s, k, i \mapsto \upd{k = 0}{s[0 := 0]}\ \ \upd{0 < k}{s[0 := s[0] + (\ite{k \bmod 2 = 0}{2}{0})]}}$:
  a sum in a cell, with the summand a function of the iteration.
\item \textbf{The cell is a sum.} The postcondition reads $s'[0] = \cur(S)[0]$;
  \lemma{select\_current\_trace\_reset\_add} matches with $j := 0$, $c := 0$ and
  $G := \lambda k.\ \ite{k \bmod 2 = 0}{2}{0}$, abstracted from the summand, and yields
  $\ite{n = 0}{s_0[0]}{0 + \mathsf{sum}(\mapc{\bl{\cnt(0, 1, n)}}{k \mapsto \ite{0 < k}{G(k)}{0}})}$; the
  case $n = 0$ is $\intf(n) = 0$ (\lemma{zero\_integer}), false under $\intf(n) = N > 0$.
\item \textbf{The bound.} \lemma{sum\_cases\_below} matches the remaining atom
  $\mathsf{sum}(\ldots) \le 2N$ with $e := 2N$, $c := 2$, $d := 0$ and $P$ the parity; its conditions
  $0 \le 2N$, $2\,\intf(n) \le 2N$ and $0 \le 2N$ hold by $\intf(n) = N$. The goal is closed.
\end{enumerate}
A variant asserting \texttt{sum[0] <= N - 2} stays open with the residual $\neg(0 < N)$: the lemmas compute
the value $N - 1$ of the sum (the summand $2$ at the even positions from $2$ on) and the bound is false; the
same lemmas close \texttt{brs3} to \texttt{brs5}, and \lemma{sum\_mod\_below} in place of the third closes
\texttt{ms2} to \texttt{ms5}, whose arrays hold $i \bmod m$ (Section~\ref{sec:ms2}).

\subsection{The idea}
The scalar version of this program, which the Dafny translations of SV-COMP produce, is closed by
\lemma{current\_trace\_reset\_add} alone; the C program needs the same lemma for a cell of an array, with
the same right-hand side. This is ``name the normal form'' of Section~\ref{sec:principles:shape} at the level
of the database: the shape of the data decides which lemma matches, and the lemma for the new shape
should land in the normal form of the old, so that everything after it is shared. The bug that this test
found on the way is finding~10 of the algorithms seen from the engine: the simplifier's eta reduction had
turned the instance $\lambda z.\ a[z]$ of $G$ into the array $a$ itself, which left an application $a(k)$
that no rule about $a[k]$ matched, fixed in version 0.35.15 by restricting eta reduction to lambdas of
array type.
